%-------------------------------------------------------------------------------
% Tanvir Ahmed - Curriculum Vitae
% Same layout as the original CV (centred header, ruled section titles, bold
% entry titles with right-aligned dates, bullet details, "Technologies:" lines),
% with the content brought up to date.
%
% Build: ./build-cv.sh   (pdflatex, two passes)
%-------------------------------------------------------------------------------
\documentclass[11pt,letterpaper]{article}
\usepackage[margin=1in]{geometry}
\usepackage[T1]{fontenc}
\usepackage[utf8]{inputenc}
\usepackage{titlesec}
\usepackage{enumitem}
\usepackage{multicol}
\usepackage{parskip}
\usepackage{array}
\usepackage[hidelinks]{hyperref}
\input{glyphtounicode}
\pdfgentounicode=1
\pagestyle{plain}

\titleformat{\section}{\Large\bfseries}{}{0em}{}[\vspace{-5pt}\titlerule]
\titlespacing*{\section}{0pt}{12pt}{7pt}
\setlist[itemize]{leftmargin=1.5em, itemsep=0pt, topsep=2pt, parsep=0pt}
\setlength{\parskip}{3.5pt}

% Entry head: title wraps inside its own column; the date or status stays on the
% first line, right-aligned. Every macro opens and closes its own paragraph so a
% following line can never run onto the previous one.
\newcommand{\entryhead}[2]{\par\vspace{3pt}\noindent%
  \begin{tabular*}{\textwidth}[t]{@{}>{\raggedright\arraybackslash}p{0.73\textwidth}@{\extracolsep{\fill}}r@{}}
    \textbf{#1} & #2
  \end{tabular*}\par\vspace{-1pt}}
\newcommand{\entry}[2]{\entryhead{#1}{\textbf{#2}}}
\newcommand{\entryplain}[2]{\entryhead{#1}{#2}}
\newcommand{\tech}[1]{\par\vspace{1pt}\noindent\textbf{Technologies:} #1\par}
\newcommand{\links}[1]{\par\vspace{1pt}\noindent\textbf{Links:} #1\par}
\newcommand{\gap}{\par\vspace{3pt}}

\begin{document}

%----------------------------------- HEADER ------------------------------------
\begin{center}
  {\LARGE\bfseries TANVIR AHMED}\\[5pt]
  North South University, Dhaka, Bangladesh\\[2pt]
  \href{mailto:tanvirahmed.minhaz1@gmail.com}{tanvirahmed.minhaz1@gmail.com} \,|\,
  \href{tel:+8801841449894}{+880 1841 449894} \,|\,
  \href{https://github.com/minhaz-42}{GitHub} \,|\,
  \href{https://www.linkedin.com/in/tanvir-ahmed-1a4856418/}{LinkedIn} \,|\,
  \href{https://minhaz-42.github.io/Tanvir-Ahmed/}{Website}
\end{center}

%---------------------------------- EDUCATION ----------------------------------
\section{Education}
\entry{Bachelor of Science in Computer Science and Engineering}{2022 -- Dec 2026 (expected)}
North South University, Dhaka, Bangladesh\\
Department of Electrical and Computer Engineering\\
CGPA: \textbf{3.02 / 4.00}
\gap
\entry{Higher Secondary Certificate (HSC), Science}{January 2021}
Milestone College, Dhaka\\
GPA: \textbf{5.00 / 5.00}

%----------------------------- RESEARCH INTERESTS ------------------------------
\section{Research Interests}
\begin{multicols}{2}
\begin{itemize}[topsep=0pt]
  \item Efficient Multimodal Inference
  \item Computer Vision and Video Understanding
  \item Calibration, Abstention and Uncertainty
  \item Quantum Machine Learning
  \item Explainable AI (XAI)
  \item Edge AI
  \item Agentic Systems
  \item Full-Stack System Architecture
\end{itemize}
\end{multicols}

%-------------------------------- PUBLICATIONS ---------------------------------
\section{Publications}
\entry{Closed-Book Knowledge Base Construction with Calibrated Abstention and Numeric Self-Consistency}{Accepted}
Tanvir Ahmed. \textit{LM-KBC 2026 shared task (5th edition), Automated Knowledge Base Construction (AKBC) at EMNLP 2026.}
\begin{itemize}
  \item Set-valued knowledge base construction with an open-weight model of at most 32B parameters,
    no retrieval and no fine-tuning; treats deciding \textit{when to answer} as a first-class decision
    and beats the official baseline on five of the six relations.
  \item Code: \href{https://github.com/minhaz-42/lm-kbc-2026}{github.com/minhaz-42/lm-kbc-2026}
\end{itemize}
\gap
\entry{SoundSentry: A Low-Cost Edge-AI Acoustic Event Detection and Direction-Finding System on STM32}{Under review}
Tanvir Ahmed, Sanzana Afrin Tonny, Husain Kabir Sifat. \textit{International Conference on Electrical and Computer Engineering (ICECE) 2026, Dhaka.}
\begin{itemize}
  \item STM32F103 prototype with a three-microphone array, ADC-DMA acquisition, on-device
    classification over 11 acoustic features, servo-based direction response and SD-card logging.
  \item Code: \href{https://github.com/minhaz-42/SoundSentry-Acoustic-Surveillance-Direction-Finding-System}{github.com/minhaz-42/SoundSentry-Acoustic-Surveillance-Direction-Finding-System}
\end{itemize}

%------------------------------ RESEARCH PROJECTS ------------------------------
\section{Research Projects}

\entry{Frame-Aware Multi-Model Scheduling for Efficient Video Captioning}{Senior Design Project, 2026}
Routes video captioning between a 4B and a 27B vision-language model so the large model is spent
only where the footage is hard. Lead author of four.
\begin{itemize}
  \item Frames are stitched into mosaics, encoded by a frozen CLIP ViT-B/32 and scored for visual
    complexity by a small graph attention network trained without human labels.
  \item Answers a 19-clip benchmark in fewer requests than either model alone while scoring above the
    small model on every judged axis.
\end{itemize}
\tech{Python, PyTorch, PyTorch Geometric, CLIP, Django, Gemma 3, Qwen3-VL}
\gap

\entry{Quantum-Enhanced Driver Distraction Detection}{Directed Research, 2026}
A hybrid classical--quantum pipeline with uncertainty-aware risk scoring, tested under a protocol
where only the classification head differs.
\begin{itemize}
  \item The variational quantum head reaches accuracy parity with its matched classical control and
    cuts expected calibration error by 37\% relative.
  \item The report states what the evidence does not support, including that the calibration gain
    cannot be attributed to the circuit.
\end{itemize}
\tech{Python, PyTorch, PennyLane, EfficientNet, scikit-learn}
\gap

\entry{Bangla Mental Health Conversational Agent}{Ongoing}
AI-powered mental health chatbot for Bangla-speaking users providing context-aware, empathetic
conversations.
\begin{itemize}
  \item Fine-tuned models: Gemma-3-Bangla, Qwen-2.5-0.5B-Bangla
  \item Features: voice I/O, mood tracking, conversation history, multi-language UI
  \item Technical: Django backend, OAuth authentication, Apple MPS GPU acceleration
\end{itemize}
\tech{Python, Django, PyTorch, Transformers, REST APIs, SQLite}
\gap

\entry{Observing Is Not Free: the decision value of visual observations in vision-language agents}{In preparation}
Measures whether an agent's observation-selection criterion ranks observations the way their
realised effect on the answer does, using a counterfactual table of answers before and after each
candidate observation.
\tech{Python, Vision-Language Models}
\gap

\entry{Knowing What You No Longer Know: calibrated belief validity for persistent world memory}{In preparation}
Separates whether a belief is still true from how confident the model was when it formed it, and
builds a benchmark of elapsed-time-matched pairs to score it.
\tech{Python, PyTorch, HD-EPIC, AI2-THOR, ProcTHOR}
\gap

\entry{SimCert: A Classical-Simulability Audit of Trained QNLP Text Models}{Abstract submitted, QTML 2026}
Re-simulates trained quantum-NLP circuits under matrix-product-state truncation to test whether the
quantum component is load-bearing. Code: \href{https://github.com/minhaz-42/simcert-qnlp}{github.com/minhaz-42/simcert-qnlp}
\tech{Python, PennyLane, lambeq, OpenQASM, Hydra}
\gap

\entry{Sequential Video Belief}{In progress}
Whether a video-language model's own stated answer impairs its ability to revise that answer when
later visual evidence contradicts it.
\tech{Python, Vision-Language Models}
\gap

\entry{Uncertainty-Guided Language-Grounded 3D Scene Graphs from Sparse Views}{In progress}
Frozen geometry and vision-language models with a lightweight graph reasoning module, for 3D
object segmentation and spatial relations from a few casual images.
\tech{Python, PyTorch, VGGT, SAM 2, CLIP, DINOv2}

%------------------------------ TECHNICAL PROJECTS -----------------------------
\section{Technical Projects}

\entryplain{Pakna Council -- AI Council Platform}{2026}
Local-first web app running multiple LLMs in parallel on one prompt through Ollama, with model
comparison and vision workflows.
\links{\href{https://github.com/minhaz-42/Pakna_council}{Source} $\cdot$ \href{https://pakna-council.vercel.app/}{Live site}}
\tech{Python, Django, React, Docker, Ollama, WebSocket}
\gap

\entryplain{SHAND -- Explainable Assumption Detector}{2026}
Rule-based system for analysing assumptions in English text; includes visual analytics and export.
\links{\href{https://github.com/minhaz-42/SHAND}{Source} $\cdot$ \href{https://shand1.vercel.app/}{Live site}}
\tech{Python, NLP, HTML, JavaScript, D3.js}
\gap

\entryplain{ISKA -- Privacy-First Document Analysis Assistant}{2026}
Document analysis assistant with a Chrome extension for use while reading.
\links{\href{https://github.com/minhaz-42/ISKA}{Source} $\cdot$ \href{https://iska-hxrr.vercel.app/}{Live site}}
\tech{TypeScript, Chrome extension APIs}
\gap

\entryplain{Nestledown -- 3D Exploration Game}{2026}
Cozy 3D exploration game built solo in Godot 4 with every asset script-generated; a parallel build
runs on a custom C++20 engine with a Metal renderer.
\links{\href{https://youtu.be/MUN7_kuh00U}{Announcement trailer}}
\tech{Godot 4, Blender, C++20, Metal, SDL3, CMake}
\gap

\entryplain{Entangled Fusion Quantum Circuits for Vision--Language Alignment}{2026}
Experimental multimodal platform combining vision analysis with quantum-inspired fusion and
natural language reasoning.
\links{\href{https://github.com/minhaz-42/Entangled-Fusion-Quantum-Circuits-for-Vision-Language-Representation-Alignment}{Source}}
\tech{Python, Django, PennyLane, Ollama, OpenCV, REST APIs}
\gap

\entryplain{SUMBA -- 3D Sign Language Translation}{2026}
Real-time sign-to-text translation using computer vision with mobile integration and offline
functionality.
\links{\href{https://github.com/minhaz-42/SUMBA}{Source}}
\tech{Django, MediaPipe, Transformers, React Native}
\gap

\entryplain{FitWell -- AI Nutrition Assistant}{2025}
AI-powered meal planning, BMI tracking and health insights.
\links{\href{https://github.com/minhaz-42/FitWell-Using-Django}{Source}}
\tech{Django, Python, Qwen, SQLite, Chart.js}
\gap

\entryplain{DG Healthcare Management System}{2024}
Hospital management platform with appointments, billing and role-based access.
\links{\href{https://github.com/minhaz-42/cse311-project-on---DG-health-care-management-system-using-html-css-php}{Source}}
\tech{HTML5, CSS3, PHP, MySQL, JavaScript}
\gap

\entryplain{CSE215 Java Application}{2024}
OOP-focused Java application with modular design, MVC pattern and unit testing.
\links{\href{https://github.com/minhaz-42/cse215-project-using-java}{Source}}
\tech{Java, JavaFX, MySQL, JUnit}

%------------------------------- TECHNICAL SKILLS ------------------------------
\section{Technical Skills}
\textbf{Languages:} Python, Java, JavaScript, TypeScript, PHP, C/C++, SQL, HTML, CSS, \LaTeX

\textbf{Machine Learning \& AI:} Deep Learning, Computer Vision, NLP, LLMs, Transformers,
Vision-Language Models, MediaPipe, Quantum ML, Fine-tuning, Calibration

\textbf{Frameworks \& Libraries:} PyTorch, PyTorch Geometric, TensorFlow, scikit-learn, Django,
React, Node.js, JavaFX, PennyLane, lambeq, OpenCV, CLIP, Hydra

\textbf{Databases:} MySQL, SQLite, PostgreSQL, MongoDB

\textbf{Software Engineering:} OOP, Design Patterns, Agile/Scrum, REST APIs, WebSocket, Docker,
CI/CD, Unit Testing

\textbf{Embedded \& Graphics:} STM32, Embedded C, ADC-DMA acquisition, Godot 4, Blender, Metal,
SDL3, CMake

\textbf{Tools \& Platforms:} Git, VS Code, Jupyter, \LaTeX, Linux, Ollama, Apple MPS

\textbf{Specializations:} Full-Stack Development, System Architecture, LLM Orchestration,
Edge AI Deployment

%--------------------------- EXTRACURRICULAR ACTIVITIES ------------------------
\section{Extracurricular Activities}
\entry{IEEE NSU Student Branch, General Member}{2023 -- Present}
\begin{itemize}
  \item Participated in technical workshops, seminars, and community outreach programs
\end{itemize}
\gap
\entry{NSU Trimodal AI Lab, Research Member}{2025 -- Present}
\begin{itemize}
  \item Engaged in research discussions and collaborative experimentation on multimodal and
    agentic AI systems
  \item Contributed to idea validation, prototyping, and literature exploration in emerging AI domains
\end{itemize}

%---------------------------------- REFERENCES ---------------------------------
\section{References}
\textbf{Dr. Md. Adnan Arefeen}\\
Assistant Professor, Dept. of ECE, North South University\\
\href{mailto:adnan.arefeen@northsouth.edu}{adnan.arefeen@northsouth.edu}

\end{document}
